\section{Related Work}

\label{sec:related}

\subsection{Intelligent Tutoring Systems}

Intelligent Tutoring Systems (ITS) have demonstrated significant learning gains across STEM domains \citep{vanlehn2011relative}.
Bayesian Knowledge Tracing (BKT), introduced by \citet{corbett1994knowledge}, models learner mastery as a hidden Markov process over knowledge components.
Deep Knowledge Tracing (DKT) extends this approach using recurrent neural networks \citep{piech2015deep}, achieving improved prediction accuracy at the cost of interpretability.
\citet{zhang2017dynamic} proposed Dynamic Key-Value Memory Networks that maintain explicit representations of knowledge concepts while leveraging neural attention mechanisms.

More recently, transformer-based models have been applied to learner modeling.
\citet{ghosh2020context} introduced Attentive Knowledge Tracing (AKT), which uses self-attention to capture complex temporal dependencies in learning sequences.
\citet{shin2021saint} further refined this approach with SAINT+, incorporating elapsed time and lag time features.

\subsection{Conservation Education Technology}

Technology-enhanced conservation education remains nascent compared to mainstream EdTech.
\citet{johnson2018wildlife} developed a mobile species identification app with embedded educational content, but without adaptive sequencing.
\citet{rousell2019gamification} surveyed gamification approaches in environmental education, finding positive engagement effects but limited evidence of knowledge transfer to field practice.
The iNaturalist platform \citep{horn2018inaturalist} combines citizen science with implicit learning through species identification feedback, though it lacks structured pedagogical progression.

\subsection{Multi-Language Educational NLP}

Cross-lingual educational content generation presents unique challenges.
\citet{fan2021beyond} demonstrated that multilingual language models can generate pedagogically appropriate explanations across languages, though quality degrades for low-resource languages.
\citet{artetxe2020cross} showed that cross-lingual transfer learning can reduce the data requirements for educational NLP in new languages by 60--80\%.
The specific challenges of translating ecological terminology across linguistic and cultural boundaries have received limited attention \citep{maffi2005linguistic}.

\subsection{Decentralized Education Platforms}

Blockchain-based education systems have been proposed for credential verification \citep{grech2017blockchain} and incentive alignment \citep{sharples2016blockchain}.
However, most implementations focus on certification rather than learning processes.
The Zoo Network's Decentralized Semantic Optimization (DSO) protocol \citep{zoo2025dso} provides a foundation for collaborative knowledge curation that we extend to educational content.

% ============================================================
% 3. Conservation Competency Ontology
% ============================================================
